\subsection{A smaller common-source gate with one disk in each first-skeleton weight}

Fix $d\geq0$.  Let $\mathcal E^{[d]}$ be the full
subcategory of $\mathcal A_{\mathbb Z}$ whose objects
have the following properties for some integer $N\geq2$:
their first skeleton is strictly isomorphic over $S$
to the particular $P(N)$ of (SK.1), their coefficient
comparison holds through form degree $d$ in every
weight and simplex, and their normalized rows in form
degrees $0,\ldots,d+2$ are exact in every weight
at least $N$.  Here $P(N)$ has exactly one even
zero-form disk generator in each weight $w\geq N$.
The choice of an isomorphism or of $N$ is a property,
not extra morphism data.  In particular
$\mathcal E^{[d]}\subset\mathcal C^{[d]}$ is full.

\begin{theorem}[A-S09-27: Proved; a strictly smaller internal common-source gate]
\label{AS0927}
The following three assertions are equivalent:
\begin{enumerate}
\item Every pair of objects in $\mathcal A_{\mathbb Z}$
has a common admissible source.
\item Every pair of objects in $\mathcal E^{[d]}$
has a common source internal to $\mathcal E^{[d]}$.
\item Every pair in $\mathcal E^{[d]}$ whose first
skeletons are both the same $P(N)$ has a common source
internal to $\mathcal E^{[d]}$.
\end{enumerate}
The inclusion $\mathcal E^{[d]}\subset\mathcal C^{[d]}$
is proper, even up to strict isomorphism.
All assertions hold with the support condition $q\leq w$,
with termwise flatness over $\mathbb Z$, and with both.
This theorem concerns the common-source assertion;
the terminal-localization gate remains A-S09-24.
\end{theorem}

\begin{proof}
We first construct actual replacements for an arbitrary
finite list of admissible objects $A_i$.  Choose one
integer
\[
 N\geq\max_i\{2,\lambda_{A_i}(d+2),
                         \mu_{A_i}(d+2)\}.
 \tag{S09.139}
\]
Use the strict tail functor $T_N$ of A-S09-24.
For each $i$, its rows through degree $d+2$ are
exact in weights at least $N$, because those rows
are literally the corresponding rows of $A_i$.
Its coefficient comparison through degree $d+2$
holds in all weights: below $N$ it is the coefficient
complex $S^w[0]$, and above $N$ it is the given
comparison for $A_i$.

For every $w\geq N$, choose one normalized coefficient
edge in $(T_NA_i)_1^{w,0}$ with endpoints $0,T^{[w]}$.
Row-zero exactness supplies these elements, so (SK.7)
gives actual strict truncated PD maps
\[
                      p_i:P(N)\longrightarrow (T_NA_i)_{\leq1}.
 \tag{S09.140}
\]
Use precisely the pullback of A-004:
\[
 E_i=T_NA_i\mathop{\times}_{C((T_NA_i)_{\leq1})}C(P(N)).
 \tag{S09.141}
\]
The strict first-skeleton identifications and the two
projections are those explicitly proved in (SK.11).
They give the actual comparison map
\[
             E_i\longrightarrow T_NA_i
                          \xrightarrow{\ t_{N,A_i}\ }A_i.
 \tag{S09.142}
\]

A-004 proves that $E_i$ is admissible.  Its horizontal
bound in form degrees through $d+2$ is $N$, by (SK.4)
and the bound for $T_NA_i$ just established.
The same theorem gives its coefficient comparison
through degree $d$ in every weight at least $N$:
both degree-$d+2$ bounds used there are at most $N$.
In weights below $N$, all objects in the pullback
(S09.141) are literally the coefficient pieces.
For $C(P(N))$ this follows from the absence of any
generator of weight less than $N$; its matching
families have all vertices and edges equal to the
same coefficient.  For $T_NA_i$ and its first-skeleton
coskeleton it follows from (S09.116) and the same
matching formula.  Their maps on these pieces are
identities, so their pullback is also $S^w[0]$.
Thus $E_i$ has the coefficient comparison through
degree $d$ in every weight, and
\[
                         E_i\in\mathcal E^{[d]}.
 \tag{S09.143}
\]
Every $E_i$ in this finite list has the same first
skeleton $P(N)$.  The support and flatness properties
follow from the actual tail pieces and the subgroup
argument in A-004, and hold in every weight.

Assertion 2 immediately implies assertion 3.
For assertion 3 implying assertion 1, apply
(S09.139)--(S09.143) to an arbitrary pair $A,B$.
Its two replacements have the same $P(N)$.
A common source $D\in\mathcal E^{[d]}$ given by
assertion 3, composed with the actual maps (S09.142),
is an admissible common source for $A,B$.

For assertion 1 implying assertion 2, take a pair
$E,F\in\mathcal E^{[d]}$ and an admissible common
source $D\to E,F$ supplied by assertion 1.
Apply the same replacement construction to the
single object $D$, with its own sufficiently large
integer $M$ in (S09.139).  It supplies an actual
map $D'\to D$ with $D'\in\mathcal E^{[d]}$.
The composites $D'\to D\to E,F$ are morphisms
of the full subcategory and are the required
internal common source.  This also explains why
no fixed value of $N$ is imposed on the source in
assertion 3.  All three implications preserve the
optional support and flatness restrictions.

We prove properness by an explicit object and a
first-skeleton invariant.  Let $G$ be the standard
integral theory and choose
\[
 N\geq\max\{2,\lambda_G(d+2),\mu_G(d+2)\}.
 \tag{S09.144}
\]
The functorial object $C_*=R(T_NG)$ of A-006 and
A-S09-24 belongs to $\mathcal C^{[d]}$ and has
the support and flatness properties.
In $G_1^0=S\otimes\Gamma(x)$ write $d_1x=0$,
$d_0x=T$.  For every integer $k$ the element
\[
 a_k=\gamma_N(x)+k\,T^{[N-2]}
                              \bigl(Tx-2\gamma_2(x)\bigr)
 \tag{S09.145}
\]
is a normalized coefficient edge of weight $N$:
its first face is zero, and its zeroth face is
$T^{[N]}$ because $T^2=2T^{[2]}$.
The elements are pairwise distinct.  In the free
PD monomial basis, the nonzero $x$-linear term
$T^{[N-2]}Tx$ and the $x$-quadratic term
$2T^{[N-2]}\gamma_2(x)$ have different monomials.
Both terms retain their nonzero integral coefficients.
Thus the filler set $E_N(T_NG)$ is infinite.
There are no positive-weight fillers below $N$,
as proved explicitly in A-S09-24.

Consequently the first positive weight in which
the quotient
\[
             (C_*)_1^{w,0}/S^w
 \tag{S09.146}
\]
is nonzero is $w=N$, and its group there is free
abelian of infinite rank: its independent basis
contains one variable for every $a_k$ in (S09.145).
There are no products of lower positive-weight
free variables, since there are no such variables.
In contrast, for the specified $P(M)$ the analogous
quotient first becomes nonzero in weight $M$, and
in that weight is free of rank one.
A coefficient-preserving graded isomorphism would
identify these quotient groups in every weight.
It would first force $M=N$ and then identify an
infinite-rank free group with a rank-one free group,
which is impossible, for example after tensoring
with $\mathbb Q$ and comparing vector-space dimensions.
Thus the first skeleton of $C_*$ is isomorphic to
no $P(M)$, and $C_*\notin\mathcal E^{[d]}$.
This proves strict properness in every stated variant.
\end{proof}

The common-source construction in this smaller category
has not been supplied.  Its exact finite-degree conditions
and both directions of the reduction are the content of
the theorem.
